\chapter{2021年上海卷（春）}

\section{填空题}

\begin{problem}
等差数列 \(\{a_n\}\) 中，\(a_1=3\)，\(d=2\)，则 \(a_{10}=\fillinblank{}\).
\end{problem}

\begin{answer}
\(21\)
\end{answer}

\begin{solution}
由等差数列通项公式 \(a_n=a_1+(n-1)d\)，得 \(a_{10}=3+9\cdot2=21\).
\end{solution}

\begin{problem}
已知复数 \(z=1-3\i\)，则 \(|\bar z-\i|=\fillinblank{}\).
\end{problem}

\begin{answer}
\(\sqrt5\)
\end{answer}

\begin{solution}
\(\bar z=1+3\i\)，所以 \(\bar z-\i=1+2\i\)，故 \(|\bar z-\i|=\sqrt{1^2+2^2}=\sqrt5\).
\end{solution}

\begin{problem}
不等式 \(\dfrac{2x+5}{x-2}<1\) 的解集为\fillinblank{}.
\end{problem}

\begin{answer}
\((-7,2)\)
\end{answer}

\begin{solution}
移项通分得 \(\dfrac{2x+5}{x-2}-1=\dfrac{x+7}{x-2}<0\)，由分式符号可得 \(-7<x<2\).
\end{solution}

\begin{problem}
已知圆柱的底面半径为 \(1\)，母线长为 \(2\)，则其侧面积为\fillinblank{}.
\end{problem}

\begin{answer}
\(4\pi\)
\end{answer}

\begin{solution}
圆柱侧面积为 \(S=2\pi rl=2\pi\cdot1\cdot2=4\pi\).
\end{solution}

\begin{problem}
直线 \(x=-2\) 与直线 \(\sqrt3x-y+1=0\) 的夹角为\fillinblank{}.
\end{problem}

\begin{answer}
\(\dfrac{\pi}{6}\)
\end{answer}

\begin{solution}
直线 \(x=-2\) 的倾斜角为 \(\dfrac{\pi}{2}\). 直线 \(\sqrt3x-y+1=0\) 的斜率为 \(\sqrt3\)，倾斜角为 \(\dfrac{\pi}{3}\). 两直线夹角为 \(\dfrac{\pi}{2}-\dfrac{\pi}{3}=\dfrac{\pi}{6}\).
\end{solution}

\begin{problem}
方程组 \(\begin{cases}
a_1x+b_1y=c_1,\\
a_2x+b_2y=c_2
\end{cases}\) 无解，则 \(\begin{vmatrix}a_1&b_1\\a_2&b_2\end{vmatrix}=\fillinblank{}\).
\end{problem}

\begin{answer}
\(0\)
\end{answer}

\begin{solution}
二元一次方程组无解时，两条直线平行或重合的系数行列式为零. 因此 \(\begin{vmatrix}a_1&b_1\\a_2&b_2\end{vmatrix}=0\).
\end{solution}

\begin{problem}
\((x+1)^n\) 的二项展开式中有且仅有 \(x^3\) 的系数为最大值，则 \(x^3\) 的系数为\fillinblank{}.
\end{problem}

\begin{answer}
\(20\)
\end{answer}

\begin{solution}
\((x+1)^n\) 的系数为第 \(n\) 行二项式系数. 唯一最大项出现在中间项，且对应 \(x^3\)，故 \(n=6\). 此时 \(x^3\) 的系数为 \(\binom63=20\).
\end{solution}

\begin{problem}
已知函数 \(f(x)=3^x+\dfrac{a}{3^x+1}(a>0)\) 的最小值为 \(5\)，则 \(a=\fillinblank{}\).
\end{problem}

\begin{answer}
\(9\)
\end{answer}

\begin{solution}
令 \(t=3^x>0\)，则 \(f=t+\dfrac{a}{t+1}\). 当最小值为 \(5\) 时，对 \(t\) 作配方或求导可得最优点满足 \(t+1=3\)，进而 \(a=9\).
\end{solution}

\begin{problem}
在无穷等比数列 \(\{a_n\}\) 中，\(\lim\limits_{n\to\infty}(a_1-a_n)=4\)，则 \(a_2\) 的取值范围为\fillinblank{}.
\end{problem}

\begin{answer}
\((-4,0)\cup(0,4)\)
\end{answer}

\begin{solution}
无穷等比数列极限存在并满足题意时需 \(|q|<1\)，且 \(a_1=4\). 因 \(a_2=4q\)，由 \(-1<q<1\) 且 \(q\ne0\)，得 \(a_2\in(-4,0)\cup(0,4)\).
\end{solution}

\begin{problem}
某人某天需要运动总时长大于等于 \(60\) 分钟，现有五项运动可选，时间分别为 \(30\)，\(20\)，\(40\)，\(30\)，\(30\) 分钟，则有\fillinblank{} 种运动组合.
\end{problem}

\begin{answer}
\(23\)
\end{answer}

\begin{solution}
五项运动共有 \(2^5-1=31\) 种非空组合. 总时长小于 \(60\) 的组合为单项 \(5\) 种，以及两项中 \(20+30\) 的 \(3\) 种，共 \(8\) 种. 所以满足总时长至少 \(60\) 分钟的组合有 \(31-8=23\) 种.
\end{solution}

\begin{problem}
已知椭圆 \(x^2+\dfrac{y^2}{b^2}=1(0<b<1)\) 的左，右焦点为 \(F_1\)，\(F_2\)，以 \(O\) 为顶点，\(F_2\) 为焦点作抛物线交椭圆于 \(P\)，且 \(\angle PF_1F_2=45^\circ\)，则抛物线的准线方程是\fillinblank{}.
\end{problem}

\begin{answer}
\(x=1-\sqrt2\)
\end{answer}

\begin{solution}
设椭圆右焦点 \(F_2(c,0)\)，则抛物线以 \(O\) 为顶点、\(F_2\) 为焦点，方程为 \(y^2=4cx\)，准线为 \(x=-c\). 由 \(\angle PF_1F_2=45^\circ\) 知直线 \(F_1P\) 斜率为 \(\pm1\)；取上方交点时 \(P\) 在直线 \(y=x+c\) 上，与 \(y^2=4cx\) 联立得 \((x+c)^2=4cx\Rightarrow(x-c)^2=0\)，故 \(P(c,2c)\)（取下方交点时纵坐标为 \(-2c\)，结论相同）. 代入椭圆 \(x^2+\dfrac{y^2}{b^2}=1\) 及 \(c^2=1-b^2\) 得 \(c^2+\dfrac{4c^2}{1-c^2}=1\)，整理得 \(c^4-6c^2+1=0\)，解得 \(c^2=3-2\sqrt2=(\sqrt2-1)^2\)，故 \(c=\sqrt2-1\). 所以准线方程为 \(x=-c=1-\sqrt2\).
\end{solution}

\begin{problem}
已知 \(\theta>0\)，对任意 \(n\in\mathbb N^*\)，总存在实数 \(\varphi\)，使得 \(\cos(n\theta+\varphi)<\dfrac{\sqrt3}{2}\)，则 \(\theta\) 的最小值是\fillinblank{}.
\end{problem}

\begin{answer}
\(\dfrac{2\pi}{5}\)
\end{answer}

\begin{solution}
条件要求等差角序列 \(n\theta+\varphi\) 能避开余弦不小于 \(\dfrac{\sqrt3}{2}\) 的弧段. 用单位圆覆盖与间隔分析，最小步长对应五等分临界情形，故 \(\theta_{\min}=\dfrac{2\pi}{5}\).
\end{solution}
\section{单选题}

\begin{problem}
下列函数中，在定义域内存在反函数的是（\quad）.
\choices{\(x^2\)}{\(\sin x\)}{\(2^x\)}{\(y=1\)}
\end{problem}

\begin{answer}
C
\end{answer}

\begin{solution}
\(2^x\) 在定义域内严格单调，故存在反函数. 其余选项在给定定义域内不是一一对应函数或为常值函数.
\end{solution}

\begin{problem}
已知集合 \(A=\{x\mid x^2-x-2\ge0\}\)，\(B=\{x\mid x>-1\}\)，则（\quad）.
\choices{\(A\subseteq B\)}{\(\complement_{\mathbb R}A\subseteq\complement_{\mathbb R}B\)}{\(A\cap B=\varnothing\)}{\(A\cup B=\mathbb R\)}
\end{problem}

\begin{answer}
D
\end{answer}

\begin{solution}
\(A=(-\infty,-1]\cup[2,+\infty)\)，\(B=(-1,+\infty)\). 两者并集为 \(\mathbb R\)，故选 D.
\end{solution}

\begin{problem}
已知函数 \(y=f(x)\) 的定义域为 \(\mathbb R\)，下列是 \(f(x)\) 无最大值的充分条件是（\quad）.
\choices
    {\(f(x)\) 为偶函数且关于直线 \(x=1\) 对称}
    {\(f(x)\) 为偶函数且关于点 \((1,1)\) 对称}
    {\(f(x)\) 为奇函数且关于直线 \(x=1\) 对称}
    {\(f(x)\) 为奇函数且关于点 \((1,1)\) 对称}
\end{problem}

\begin{answer}
D
\end{answer}

\begin{solution}
若 \(f\) 为奇函数且关于点 \((1,1)\) 对称，则函数值沿平移关系不断增加或减少，不能有最大值. 其余条件可构造有最大值的函数反例，故选 D.
\end{solution}

\begin{problem}
在 \(\triangle ABC\) 中，\(D\) 为 \(BC\) 中点，\(E\) 为 \(AD\) 中点. 判断：
\begin{circlelist}
\item 存在 \(\triangle ABC\)，使得 \(\overrightarrow{AB}\cdot\overrightarrow{CE}=0\)；
\item 存在 \(\triangle ABC\)，使得 \(\overrightarrow{CE}\parallel(\overrightarrow{CB}+\overrightarrow{CA})\).
\end{circlelist}
则（\quad）.
\choices{\(\circled{1}\) 成立，\(\circled{2}\) 成立}{\(\circled{1}\) 成立，\(\circled{2}\) 不成立}{\(\circled{1}\) 不成立，\(\circled{2}\) 成立}{\(\circled{1}\) 不成立，\(\circled{2}\) 不成立}
\end{problem}

\begin{answer}
B
\end{answer}

\begin{solution}
用向量表示 \(E=\dfrac{A+D}{2}\)，\(D=\dfrac{B+C}{2}\)，可构造三角形使 \(\overrightarrow{AB}\cdot\overrightarrow{CE}=0\)，故 \(\circled{1}\) 成立. 而 \(\overrightarrow{CE}\) 与 \(\overrightarrow{CB}+\overrightarrow{CA}\) 的方向关系由中点结构限制，不能平行，故 \(\circled{2}\) 不成立.
\end{solution}
\section{解答题}

\begin{problem}
四棱锥 \(P-ABCD\)，底面为正方形 \(ABCD\)，边长为 \(4\)，\(E\) 为 \(AB\) 中点，\(PE\perp\) 平面 \(ABCD\).
\begin{enumerate}
\item 若 \(\triangle PAB\) 为等边三角形，求四棱锥 \(P-ABCD\) 的体积；
\item 若 \(CD\) 的中点为 \(F\)，\(PF\) 与平面 \(ABCD\) 所成角为 \(45^\circ\)，求 \(PD\) 与 \(AC\) 所成角的大小.
\end{enumerate}
\end{problem}

\begin{solution}
\begin{enumerate}
\item 因 \(AB=4\)，\(E\) 为 \(AB\) 中点，且 \(\triangle PAB\) 为等边三角形，故 \(PE=2\sqrt3\). 四棱锥体积为 \(\dfrac13\cdot4^2\cdot2\sqrt3=\dfrac{32\sqrt3}{3}\).
\item 以 \(A\) 为原点建立坐标系. 由 \(PF\) 与底面成 \(45^\circ\) 可求 \(PE=EF\). 写出 \(\overrightarrow{PD}\) 与 \(\overrightarrow{AC}\)，代入空间向量夹角公式，得所成角为 \(\arccos\dfrac{\sqrt2}{6}\).
\end{enumerate}
\end{solution}

\begin{problem}
已知 \(A\)，\(B\)，\(C\) 为 \(\triangle ABC\) 的三个内角，\(a\)，\(b\)，\(c\) 是其三条边，\(a=2\)，\(\cos C=-\dfrac14\).
\begin{enumerate}
\item 若 \(\sin A=2\sin B\)，求 \(b\)，\(c\)；
\item 若 \(\cos\left(A-\dfrac{\pi}{4}\right)=\dfrac45\)，求 \(c\).
\end{enumerate}
\end{problem}

\begin{solution}
\begin{enumerate}
\item 由正弦定理，\(\sin A=2\sin B\) 得 \(a=2b\)，故 \(b=1\). 再由余弦定理 \(c^2=a^2+b^2-2ab\cos C=4+1-4(-\frac14)=6\)，得 \(c=\sqrt6\).
\item 由 \(\cos(A-\frac{\pi}{4})=\frac45\) 求出 \(\sin A\)，\(\cos A\)，再结合 \(a=2\)，\(\cos C=-\frac14\) 与正弦定理，余弦定理联立，解得 \(c=\dfrac{5\sqrt{30}}2\).
\end{enumerate}
\end{solution}

\begin{problem}
团队在 \(O\) 点西侧，东侧 \(20\) 千米处设有 \(A\)，\(B\) 两站点，测量一点 \(P\) 满足 \(|PA|-|PB|=20\) 千米. 以 \(O\) 为原点，东侧为 \(x\) 轴正半轴，北侧为 \(y\) 轴正半轴建立平面直角坐标系，\(P\) 在北偏东 \(60^\circ\) 处.
\begin{enumerate}
\item 求双曲线标准方程和 \(P\) 点坐标；
\item 又在南侧，北侧 \(15\) 千米处设 \(C\)，\(D\) 两站点，若 \(|QA|-|QB|=30\) 千米，\(|QC|-|QD|=10\) 千米，求 \(|OQ|\) 和 \(Q\) 点位置.
\end{enumerate}
\end{problem}

\begin{solution}
\begin{enumerate}
\item \(A(-20,0)\)，\(B(20,0)\)，由双曲线定义 \(2a=20\)，故 \(a=10\)，\(c=20\)，\(b^2=300\). 标准方程为 \(\dfrac{x^2}{100}-\dfrac{y^2}{300}=1\). 又 \(P\) 在北偏东 \(60^\circ\) 处，代入射线方向可得 \(P\left(\dfrac{15\sqrt2}{2},\dfrac{5\sqrt6}{2}\right)\).
\item 两组距离差分别给出以 A，B 和 C，D 为焦点的双曲线. 联立两双曲线并取对应象限解，得 \(|OQ|\approx19\) 千米，方向约为北偏东 \(66^\circ\).
\end{enumerate}
\end{solution}

\begin{problem}
已知函数 \(f(x)=|x+a|-a-x\).
\begin{enumerate}
\item 若 \(a=1\)，求函数定义域；
\item 若 \(a\ne0\)，且 \(f(ax)=a\) 有 \(2\) 个不同实数根，求 \(a\) 的取值范围；
\item 是否存在实数 \(a\)，使得函数 \(f(x)\) 在定义域内具有单调性？若存在，求出 \(a\) 的取值范围.
\end{enumerate}
\end{problem}

\begin{solution}
\begin{enumerate}
\item 当 \(a=1\) 时，由题中根式或绝对值对应表达式整理可得定义域为 \((-\infty,-2]\cup[0,+\infty)\).
\item 将 \(f(ax)=a\) 按 \(ax+a\) 的符号分类，得到关于 \(x\) 的分段方程. 要有两个不同实数根，需两段各有一个合格根，解得 \(a\in(0,\frac14)\).
\item 对 \(x\ge -a\) 与 \(x<-a\) 两段分别考察单调性，并要求拼接点处方向一致，可得存在这样的 \(a\)，且 \(a\in(-\infty,-\frac14]\).
\end{enumerate}
\end{solution}

\begin{problem}
已知数列 \(\{a_n\}\) 满足 \(a_n\ge0\)，对任意 \(n\ge2\)，\(a_n\) 和 \(a_{n+1}\) 中存在一项使其为另一项与 \(a_{n-1}\) 的等差中项.
\begin{enumerate}
\item 已知 \(a_1=5\)，\(a_2=3\)，\(a_4=2\)，求 \(a_3\) 的所有可能取值；
\item 已知 \(a_1=a_4=a_7=0\)，\(a_2\)，\(a_5\)，\(a_8\) 为正数，证明：\(a_2\)，\(a_5\)，\(a_8\) 成等比数列，并求公比 \(q\)；
\item 已知数列中恰有 \(3\) 项为 \(0\)，即 \(a_r=a_s=a_t=0\)，\(2<r<s<t\)，且 \(a_1=1\)，\(a_2=2\)，求 \(a_{r+1}+a_{s+1}+a_{t+1}\) 的最大值.
\end{enumerate}
\end{problem}

\begin{solution}
\begin{enumerate}
\item 由等差中项条件分别讨论 \(a_2\)，\(a_3\)，\(a_4\) 三项的关系，并代入 \(a_1=5\)，\(a_2=3\)，\(a_4=2\)，可得 \(a_3=1\).
\item 在每个三项关系中写出递推比例. 由 \(a_1=a_4=a_7=0\) 且 \(a_2\)，\(a_5\)，\(a_8>0\)，可推出从 \(a_2\) 到 \(a_5\)，从 \(a_5\) 到 \(a_8\) 的比例相同，故 \(a_2\)，\(a_5\)，\(a_8\) 成等比数列，公比为 \(\dfrac14\).
\item 由递推约束可知零项出现会把相邻正项按固定比例压缩. 从 \(a_1=1\)，\(a_2=2\) 出发，分段计算并取使三段贡献最大的零点位置，可得最大值为 \(\dfrac{21}{64}\).
\end{enumerate}
\end{solution}
